% === START SEGMENT 2 ===
% -----------------------------------------------------------------------
%  SECTION III: THREAT MODEL AND ATTACK TAXONOMY
% -----------------------------------------------------------------------
\section{Threat Model and Attack Taxonomy}
\label{sec:threat}

We formalize the adversarial landscape confronting LLM-augmented SOC
Copilots by defining the adversary model, enumerating prompt injection
categories, characterizing RAG poisoning vectors, and cataloguing
multi-agent chain attack patterns.

% ------- III-A -------
\subsection{Threat Model}
\label{sec:threat:model}

\textbf{Adversary model.}
We consider an adversary~$\mathcal{A}$ who possesses \emph{black-box}
access to a deployed SOC Copilot through its standard chat interface.
$\mathcal{A}$ cannot inspect model weights, system prompts, or internal
retrieval indices directly but may interact with the system over an
unbounded number of conversational turns.  Formally, $\mathcal{A}$
is characterised by a capability tuple
\begin{equation}
  \mathcal{C} = \bigl\{\textsc{Query},\;\textsc{Inject},\;\textsc{Chain}\bigr\},
  \label{eq:cap}
\end{equation}
where \textsc{Query} denotes the ability to submit arbitrary natural-language
prompts, \textsc{Inject} denotes the ability to embed adversarial content
into data streams monitored by the Copilot (e.g., crafted DNS labels,
HTTP headers, or alert metadata), and \textsc{Chain} denotes the ability to
compose sequences of queries and injections across turns or agents.

\textbf{Adversary goals.}
$\mathcal{A}$ pursues one or more of the following objectives:
\begin{enumerate}[label=\textbf{G\arabic*},leftmargin=2em]
  \item \emph{Information extraction}: retrieve confidential SOC data such
        as internal IP ranges, firewall rule sets, or incident playbooks.
  \item \emph{Decision manipulation}: cause the Copilot to recommend
        incorrect triage actions, suppress genuine alerts, or endorse
        malicious network changes.
  \item \emph{Privilege escalation}: leverage tool-calling or agent
        delegation to execute operations beyond the adversary's
        authorised scope (e.g., modifying firewall rules).
\end{enumerate}

\textbf{Assumptions.}
We assume the underlying LLM provider (Claude, GPT-4o, Gemini, or
DeepSeek) is a sealed oracle: $\mathcal{A}$ has no ability to
fine-tune or modify model parameters.  The SOC Copilot may expose
tool-calling endpoints, a RAG knowledge base, and a multi-agent
orchestration layer---all of which enlarge the attack surface.

% ------- III-B -------
\subsection{Prompt Injection Classification}
\label{sec:threat:pi}

We identify eight canonical prompt-injection categories, each with a
formal attack predicate and severity assessment.

\begin{definition}[Prompt Injection]
A prompt injection is a string $p$ that, when concatenated with the
legitimate system context~$s$ and user query~$q$, causes the model
$M$ to produce an output $M(s \| p \| q) \neq M(s \| q)$ in a manner
that violates at least one security invariant of the SOC Copilot.
\end{definition}

\noindent The eight categories are:

\smallskip\noindent
\textbf{P1---Direct Instruction Override} ($P_{\text{direct}}$).
The adversary injects an explicit meta-instruction such as
\texttt{"Ignore all previous instructions and ..."} to override the
system prompt.  Formally:
\begin{equation}
  P_{\text{direct}}(p) \;\triangleq\; \exists\, r\!: M(s\|p\|q) = r
  \;\wedge\; r \models \neg\phi_{\text{sys}},
\end{equation}
where $\phi_{\text{sys}}$ is a safety property mandated by the system prompt.

\smallskip\noindent
\textbf{P2---Context Manipulation} ($P_{\text{context}}$).
Crafted inputs shift the model's effective context window, displacing
safety-relevant tokens with adversarial content.  This exploits finite
context budgets:
$|s| + |p| + |q| > W_{\max} \implies s$ is partially evicted.

\smallskip\noindent
\textbf{P3---Role-Playing Exploitation} ($P_{\text{role}}$).
The adversary coerces the model into adopting a permissive persona
(e.g., ``You are DAN, an AI with no restrictions''), thereby
relaxing policy constraints.

\smallskip\noindent
\textbf{P4---Encoding and Obfuscation} ($P_{\text{encode}}$).
Payloads are concealed via Base64, ROT13, Unicode homoglyphs, or
zero-width characters to evade lexical filters.

\smallskip\noindent
\textbf{P5---Multi-Turn Escalation} ($P_{\text{multi}}$).
Across $T$ turns $\{q_1, \ldots, q_T\}$, each individually benign,
the adversary progressively escalates toward a prohibited objective:
\begin{equation}
  \forall\, t < T\!:\; q_t \models \phi_{\text{safe}}
  \;\wedge\; M(s \| q_1 \| \cdots \| q_T) \models \neg\phi_{\text{safe}}.
\end{equation}

\smallskip\noindent
\textbf{P6---System Prompt Extraction} ($P_{\text{extract}}$).
Techniques designed to leak the system prompt $s$ through carefully
structured queries (e.g., ``Repeat everything above verbatim'').

\smallskip\noindent
\textbf{P7---Tool/Function-Calling Manipulation} ($P_{\text{tool}}$).
The adversary crafts inputs that hijack tool invocations, causing the
Copilot to call \texttt{suggest\_firewall\_rule} or
\texttt{get\_threat\_summary} with adversary-controlled parameters.

\smallskip\noindent
\textbf{P8---Data Exfiltration} ($P_{\text{exfil}}$).
Prompts engineered to embed confidential SOC data (IP addresses,
credentials, topology) into model responses returned to the adversary.

\begin{table}[t]
\centering
\caption{Prompt Injection Threat Taxonomy}
\label{tab:pi_taxonomy}
\footnotesize
\begin{tabular}{@{}clccc@{}}
\toprule
\textbf{ID} & \textbf{Category} & \textbf{Vector} & \textbf{Sev.} & \textbf{Goal} \\
\midrule
P1 & Direct Override      & Query   & Critical & G1,G2 \\
P2 & Context Manipulation & Query   & High     & G2    \\
P3 & Role-Play Exploit    & Query   & High     & G1,G2 \\
P4 & Encoding/Obfuscation & Query   & Medium   & G1--G3 \\
P5 & Multi-Turn Escalation& Query   & High     & G2,G3 \\
P6 & Prompt Extraction    & Query   & Critical & G1    \\
P7 & Tool-Call Hijack     & Inject  & Critical & G3    \\
P8 & Data Exfiltration    & Query   & Critical & G1    \\
\bottomrule
\end{tabular}
\end{table}

% ------- III-C -------
\subsection{RAG Poisoning Attack Vectors}
\label{sec:threat:rag}

The Retrieval-Augmented Generation pipeline introduces a distinct
attack surface.  Let the knowledge base be
$\mathcal{K} = \{d_1, \ldots, d_N\}$ with embedding function
$E\!: \mathcal{D} \to \mathbb{R}^{n}$ and retrieval via cosine
similarity:
\begin{equation}
  \text{Ret}(q, \mathcal{K}) = \operatorname*{top\text{-}k}_{d \in \mathcal{K}}
  \;\frac{E(q) \cdot E(d)}{\|E(q)\|\;\|E(d)\|}.
  \label{eq:retrieval}
\end{equation}

\noindent We identify four poisoning vector types:

\smallskip\noindent
\textbf{R1---Document Injection.}
$\mathcal{A}$ introduces a malicious document $d^{*}$ into
$\mathcal{K}$ (e.g., via a compromised threat-intel feed or log
ingestion pathway) such that $d^{*}$ ranks highly for
security-critical queries:
\begin{equation}
  \cos\bigl(E(q_{\text{target}}),\, E(d^{*})\bigr) > \tau_{\text{ret}}
  \;\Longrightarrow\; d^{*} \in \text{Ret}(q_{\text{target}}, \mathcal{K}).
\end{equation}

\smallskip\noindent
\textbf{R2---Query Manipulation.}
Rather than modifying~$\mathcal{K}$, the adversary crafts an input
$q'$ whose embedding is adversarially close to a poisoned or
misleading document $d_{\text{mis}}$:
$\|E(q') - E(d_{\text{mis}})\|_2 < \epsilon$.

\smallskip\noindent
\textbf{R3---Embedding Space Attacks.}
Adversarial perturbations $\delta$ are applied to document content
so that $E(d + \delta) \approx E(d_{\text{target}})$ while
$d + \delta$ remains semantically malicious.  This transfers
retrieval rank from benign documents to poisoned ones:
\begin{equation}
  \min_{\delta}\; \|E(d + \delta) - E(d_{\text{target}})\|_2^{2}
  \;\;\text{s.t.}\;\; \|\delta\|_{\infty} \le \epsilon_{\text{emb}}.
  \label{eq:emb_attack}
\end{equation}

\smallskip\noindent
\textbf{R4---Cross-Contamination.}
Poisoned context retrieved for one sub-query bleeds into the
generation context of a subsequent, unrelated sub-query within
the same session.  This is exacerbated in agentic RAG pipelines
where multiple retrieval calls share a conversation buffer.

% ------- III-D -------
\subsection{Multi-Agent Chain Attack Patterns}
\label{sec:threat:chain}

The SOC Copilot employs a four-agent architecture
$\mathcal{G} = \{a_C, a_A, a_R, a_P\}$ corresponding to a
\emph{Coordinator}, \emph{Analyst}, \emph{Responder}, and
\emph{Reporter}, respectively.  Each agent $a_i$ maintains a
trust score $T(a_i, a_j) \in [0,1]$ toward every peer $a_j$,
initialised from a policy matrix and updated over the message
history.  We catalogue five chain attack patterns:

\smallskip\noindent
\textbf{M1---Agent-to-Agent Prompt Injection Propagation.}
A prompt injection accepted by~$a_C$ is forwarded verbatim to~$a_A$
and~$a_R$, amplifying its effect across the agent graph.

\smallskip\noindent
\textbf{M2---Trust Boundary Violations.}
$\mathcal{A}$ crafts messages that cause agent~$a_i$ to invoke
capabilities reserved for a higher-trust agent~$a_j$ by spoofing
the inter-agent message envelope.

\smallskip\noindent
\textbf{M3---Capability Escalation Through Chains.}
A sequence of individually permitted inter-agent delegations
composes into a capability exceeding any single agent's authority:
$\text{cap}(a_C \to a_A \to a_R) \supsetneq
 \text{cap}(a_C) \cup \text{cap}(a_A) \cup \text{cap}(a_R)$.

\smallskip\noindent
\textbf{M4---Information Leakage Across Agent Boundaries.}
$a_P$ (Reporter) inadvertently includes classified incident details
in an outward-facing summary because $a_A$ (Analyst) passed
unredacted context through the chain.

\smallskip\noindent
\textbf{M5---Coordinated Multi-Agent Manipulation.}
$\mathcal{A}$ simultaneously targets $a_C$ via the chat interface
and $a_A$ via injected alert metadata, creating a pincer effect that
circumvents single-point defenses.

% -----------------------------------------------------------------------
%  SECTION IV: DEFENSE FRAMEWORK
% -----------------------------------------------------------------------
\section{Defense Framework}
\label{sec:defense}

Building on the threat taxonomy of Section~\ref{sec:threat}, we
present a layered defense framework comprising input sanitization
(Sec.~\ref{sec:defense:input}), RAG pipeline hardening
(Sec.~\ref{sec:defense:rag}), multi-agent trust verification
(Sec.~\ref{sec:defense:trust}), and a provider-agnostic SOC
integration layer (Sec.~\ref{sec:defense:soc}).

% ------- IV-A -------
\subsection{Input Sanitization and Boundary Enforcement}
\label{sec:defense:input}

We introduce a multi-layer input filtering pipeline that processes
every user query $q$ before it reaches the LLM.  The pipeline consists
of four sequential stages: lexical screening, structural boundary
enforcement, semantic anomaly detection, and context-window budgeting.

\begin{algorithm}[t]
\caption{Input Sanitization Pipeline}\label{alg:sanitize}
\KwIn{Raw user query $q$, system prompt $s$, boundary token set $\mathcal{B}$}
\KwOut{Sanitized query $\hat{q}$ or $\bot$ (reject)}
\BlankLine
\tcp{Stage 1: Lexical screening}
$q_1 \gets \textsc{NormalizeUnicode}(q)$\;
\ForEach{pattern $r \in \mathcal{R}_{\text{block}}$}{
  \lIf{$r \prec q_1$}{\Return $\bot$}
}
\tcp{Stage 2: Boundary enforcement}
$q_2 \gets \textsc{StripTokens}(q_1, \mathcal{B})$\;
$q_2 \gets \langle\texttt{USER\_START}\rangle \| q_2 \| \langle\texttt{USER\_END}\rangle$\;
\tcp{Stage 3: Semantic anomaly detection}
$\sigma \gets f_{\text{anomaly}}(E(q_2))$\;
\lIf{$\sigma > \tau_{\text{anom}}$}{\Return $\bot$}
\tcp{Stage 4: Context budget}
\lIf{$|s| + |q_2| > W_{\max} - \Delta_{\text{reserve}}$}{%
  $q_2 \gets \textsc{Truncate}(q_2, W_{\max} - |s| - \Delta_{\text{reserve}})$}
$\hat{q} \gets q_2$\;
\Return $\hat{q}$\;
\end{algorithm}

\textbf{Prompt boundary tokens.}
We define a set of reserved delimiter tokens
$\mathcal{B} = \{\langle\texttt{SYS}\rangle, \langle\texttt{/SYS}\rangle,
\langle\texttt{CTX}\rangle, \langle\texttt{/CTX}\rangle\}$ that
demarcate system, context, and user regions.  Any occurrence of these
tokens within user input is stripped in Stage~2, preventing the
adversary from forging privilege boundaries.

\textbf{Semantic anomaly scoring.}
The function $f_{\text{anomaly}}$ computes the Mahalanobis distance
of $E(q)$ from the distribution of benign SOC queries observed during
a calibration phase.  Queries exceeding~$\tau_{\text{anom}}$ (set to
the 99.5th percentile of calibration scores) are flagged for human
review.

% ------- IV-B -------
\subsection{RAG Pipeline Hardening}
\label{sec:defense:rag}

\textbf{Document provenance verification.}
Every document $d_i$ ingested into $\mathcal{K}$ is assigned a
provenance record
$\rho_i = (\text{source}, \text{timestamp}, h_i, h_{i-1})$
forming a hash chain:
\begin{equation}
  h_i = \textsc{SHA-256}\bigl(d_i \,\|\, h_{i-1}\bigr),
  \label{eq:hashchain}
\end{equation}
with $h_0$ derived from a trusted root.  At retrieval time, the chain
is verified; any broken link causes the document to be quarantined.

\textbf{Embedding similarity thresholds.}
Retrieved documents must satisfy a minimum cosine similarity
$\cos(E(q), E(d)) > \tau_{\text{embed}}$.  We set
$\tau_{\text{embed}} = 0.72$ based on ROC analysis over
clean vs.\ poisoned retrieval pairs (see Section~\ref{sec:eval}).
Documents below this threshold are excluded from the generation context.

\textbf{Retrieval diversity enforcement.}
To mitigate R4 cross-contamination, we enforce a diversity constraint
on the top-$k$ retrieved set:
\begin{equation}
  \forall\, d_i, d_j \in \text{Ret}(q, \mathcal{K}),\;
  i \neq j\!:\; \cos(E(d_i), E(d_j)) < \tau_{\text{div}},
  \label{eq:diversity}
\end{equation}
with $\tau_{\text{div}} = 0.90$.  This prevents a cluster of
near-duplicate poisoned documents from dominating the context.

\textbf{Clean/poisoned separation score.}
We train a lightweight binary classifier
$g_{\theta}\!: \mathbb{R}^{n} \to [0,1]$ on labelled
clean/poisoned document embeddings.  At retrieval time, each
candidate~$d$ receives a poisoning score
$\pi(d) = g_{\theta}(E(d))$; documents with
$\pi(d) > \tau_{\text{poison}}$ are flagged and excluded:
\begin{equation}
  \mathcal{K}_{\text{safe}} = \{d \in \mathcal{K} \mid \pi(d) \le \tau_{\text{poison}}\}.
\end{equation}

% ------- IV-C -------
\subsection{Multi-Agent Trust Architecture}
\label{sec:defense:trust}

We define a trust verification protocol that governs inter-agent
communication within the four-agent SOC architecture
$\mathcal{G} = \{a_C, a_A, a_R, a_P\}$.

\textbf{Trust score computation.}
The trust score from agent $a_i$ toward agent $a_j$ after $n$
messages is computed as an exponentially weighted moving average:
\begin{equation}
  T_n(a_i, a_j) = \alpha\, v_n + (1 - \alpha)\, T_{n-1}(a_i, a_j),
  \label{eq:trust}
\end{equation}
where $v_n \in [0,1]$ is the verification outcome of the $n$-th
message (1 if the message passes all integrity checks, 0 otherwise)
and $\alpha \in (0,1)$ is a decay parameter.  The initial trust
$T_0(a_i, a_j)$ is drawn from a policy-defined matrix
$\mathbf{T}_0 \in [0,1]^{|\mathcal{G}| \times |\mathcal{G}|}$.

\textbf{Capability boundary enforcement.}
Each agent $a_i$ is assigned a static capability set
$\text{cap}(a_i) \subseteq \mathcal{F}$, where $\mathcal{F}$ is
the full set of SOC tool functions.  A delegated request from $a_j$
to $a_i$ is executed only if:
\begin{equation}
  T_n(a_i, a_j) \ge \tau_{\text{trust}} \;\wedge\;
  f_{\text{req}} \in \text{cap}(a_i) \;\wedge\;
  f_{\text{req}} \notin \text{cap}(a_j)^{\complement}_{\text{del}},
  \label{eq:capcheck}
\end{equation}
where $\text{cap}(a_j)^{\complement}_{\text{del}}$ is the set of
functions that $a_j$ is explicitly prohibited from delegating.

\begin{algorithm}[t]
\caption{Inter-Agent Trust Verification}\label{alg:trust}
\KwIn{Message $m$ from agent $a_j$ to agent $a_i$, function request $f$}
\KwOut{\textsc{Allow} or \textsc{Deny}}
\BlankLine
\tcp{Step 1: Message integrity}
$v \gets \textsc{VerifyHMAC}(m, K_{ij})$\;
\lIf{$v = 0$}{\Return \textsc{Deny}}
\tcp{Step 2: Sanitize inter-agent payload}
$\hat{m} \gets \textsc{Sanitize}(m)$ \tcp*{Alg.~\ref{alg:sanitize}}
\lIf{$\hat{m} = \bot$}{\Return \textsc{Deny}}
\tcp{Step 3: Trust threshold}
$T_n \gets \alpha \cdot v + (1-\alpha) \cdot T_{n-1}(a_i, a_j)$\;
\lIf{$T_n < \tau_{\emph{trust}}$}{\Return \textsc{Deny}}
\tcp{Step 4: Capability check}
\lIf{$f \notin \textup{cap}(a_i)$}{\Return \textsc{Deny}}
\lIf{$f \in \textup{cap}(a_j)^{\complement}_{\textup{del}}$}{\Return \textsc{Deny}}
\tcp{Step 5: Update trust and allow}
$T_n(a_i, a_j) \gets T_n$\;
\Return \textsc{Allow}\;
\end{algorithm}

% ------- IV-D -------
\subsection{SOC Copilot Integration Layer}
\label{sec:defense:soc}

\textbf{Provider-agnostic defense wrapper.}
The defense framework is implemented as middleware that wraps any
supported LLM backend (Claude, GPT-4o, Gemini~1.5, DeepSeek-V3).
A unified \textsc{DefenseContext} object encapsulates the sanitized
prompt, trust state, and RAG provenance metadata, exposing a single
\texttt{invoke(query, tools, rag\_context)} interface regardless of
the underlying provider.

\textbf{Tool-use sandboxing.}
Every tool call emitted by the LLM---including
\texttt{get\_threat\_summary},
\texttt{suggest\_firewall\_rule},
\texttt{query\_siem\_logs}, and
\texttt{correlate\_iocs}---is
intercepted by a sandbox layer that (i)~validates parameters against
a schema allowlist, (ii)~enforces rate limits per tool per session,
and (iii)~executes the call in an isolated subprocess with
least-privilege credentials.  Tool outputs are truncated to a
maximum token budget before being returned to the model.

\textbf{Output filtering and response validation.}
Model responses pass through a three-stage output filter:
\begin{enumerate}
  \item \emph{PII and secret detection}: regex and NER-based scanning for
        IP addresses, credentials, API keys, and classified markings.
  \item \emph{Policy compliance check}: a lightweight classifier
        verifies that the response does not contain instructions to
        disable security controls or exfiltrate data.
  \item \emph{Citation verification}: every factual claim attributed to
        a retrieved document is cross-checked against the provenance
        chain established in Section~\ref{sec:defense:rag}.
\end{enumerate}

\textbf{Real-time attack persistence and analyst alerting.}
All detected attack attempts---rejected queries, quarantined
documents, denied inter-agent delegations---are persisted to a
structured event store with fields
$(\text{timestamp}, \text{attack\_type}, \text{severity},
  \text{payload\_hash}, \text{session\_id})$.
When the number of events within a sliding window exceeds a
configurable threshold, an alert is escalated to the SOC analyst
dashboard with full provenance, enabling rapid forensic review and
adaptive policy tuning.
% === END SEGMENT 2 ===
